\documentclass[11pt]{article}

\usepackage[margin=2.3cm]{geometry}
\usepackage{fontspec}
\usepackage{polyglossia}
\setmainlanguage{russian}
\setotherlanguage{english}
\newfontfamily\cyrillicfont{PT Serif}
\newfontfamily\cyrillicfonttt{PT Mono}
\setmainfont{PT Serif}
\setmonofont{PT Mono}

\usepackage{amsmath, amssymb}
\usepackage{enumitem}
\usepackage{booktabs}
\usepackage{xcolor}
\usepackage{titlesec}
\usepackage{parskip}

\definecolor{accent}{HTML}{2a78d6}
\definecolor{accent2}{HTML}{eb6834}

\titleformat{\section}{\large\bfseries\color{accent}}{\thesection}{0.5em}{}
\titleformat{\subsection}{\bfseries}{\thesubsection}{0.5em}{}

\setlist[itemize]{leftmargin=1.4em, topsep=2pt, itemsep=2pt}

\title{\vspace{-1.5em}Методы бутстрапа --- обзор}
\author{}
\date{}

\begin{document}
\maketitle
\vspace{-2.5em}

Справочный обзор вариантов бутстрапа --- чтобы потом самому отобрать и адаптировать под слайды семинара. Сгруппировано по тому, какую проблему решает каждый метод.

\section{Базовое деление}

\subsection*{Непараметрический (case / Efron's) бутстрап}
\begin{itemize}
  \item \textbf{Идея:} ресэмплируем $X_1^*,\dots,X_n^*$ с возвращением прямо из наблюдаемых данных $X_1,\dots,X_n$. Никаких предположений о форме $F$ --- используем $\widehat F_n$ (эмпирическую CDF) как замену $F$.
  \item \textbf{Когда использовать:} не хочется предполагать конкретное параметрическое семейство для $F$; стандартный выбор для стандартных ошибок / доверительных интервалов на статистиках вроде среднего, медианы, корреляции, коэффициентов регрессии.
  \item \textbf{На что смотреть:} нужны i.i.d.\ данные; не работает для статистик, зависящих от экстремумов (например, минимум/максимум выборки --- см.\ \S\ref{sec:fails}).
\end{itemize}

\subsection*{Параметрический бутстрап}
\begin{itemize}
  \item \textbf{Идея:} сначала подгоняем параметрическую модель $F_{\hat\theta}$ к данным (например, $\hat\theta = (\hat\mu,\hat\sigma^2)$ для $N(\mu,\sigma^2)$), затем генерируем бутстрап-выборки \emph{из подогнанной модели} $F_{\hat\theta}$, а не из сырых данных.
  \item \textbf{Когда использовать:} готовы предположить модель и хотите заложить это предположение прямо в оценку неопределённости (например, симуляция поведения пула при заданном процессе цены).
  \item \textbf{На что смотреть:} оценка неопределённости настолько хороша, насколько верно параметрическое предположение; ошибка в модели незаметно искажает всё, что построено поверх неё.
\end{itemize}

\section{Расширения для структурированных / зависимых данных}

\subsection*{Residual bootstrap (регрессия)}
\begin{itemize}
  \item \textbf{Идея:} подгоняем $Y_i = \hat Y_i + \hat\varepsilon_i$, ресэмплируем сами \emph{остатки} $\hat\varepsilon_i$ с возвращением, добавляем их обратно к предсказанным значениям, получая новые синтетические отклики $Y_i^* = \hat Y_i + \varepsilon_i^*$.
  \item \textbf{Когда использовать:} задача регрессии, ошибки примерно i.i.d.\ и гомоскедастичны.
\end{itemize}

\subsection*{Wild bootstrap}
\begin{itemize}
  \item \textbf{Идея:} как residual bootstrap, но вместо перестановки остатков каждый остаток умножается на независимый случайный вес $v_i$ (обычно $\pm1$ с вероятностью $\tfrac12$): $\varepsilon_i^* = v_i\,\hat\varepsilon_i$.
  \item \textbf{Когда использовать:} регрессия с \emph{гетероскедастичными} ошибками --- residual bootstrap здесь ломается, потому что перемешивает остатки между наблюдениями с разной дисперсией; wild bootstrap оставляет масштаб каждого остатка привязанным к своему наблюдению.
\end{itemize}

\subsection*{Block bootstrap (moving / circular block)}
\begin{itemize}
  \item \textbf{Идея:} ресэмплируются целые непрерывные \emph{блоки} последовательных наблюдений (длины $\ell$), а не отдельные точки, затем блоки склеиваются в синтетический ряд.
  \item \textbf{Когда использовать:} временные ряды / зависимые данные --- сохраняет локальную зависимость внутри блока.
  \item \textbf{На что смотреть:} результат чувствителен к длине блока $\ell$; слишком короткий $\Rightarrow$ зависимость разрушается, слишком длинный $\Rightarrow$ слишком мало блоков для ресэмплинга.
\end{itemize}

\subsection*{Stationary bootstrap}
\begin{itemize}
  \item \textbf{Идея:} block bootstrap со случайной (геометрически распределённой) длиной блока вместо фиксированной $\ell$ (Politis \& Romano, 1994).
  \item \textbf{Когда использовать:} то же, что и block bootstrap, но не хочется, чтобы результат зависел от выбора одной фиксированной длины блока.
\end{itemize}

\subsection*{Bayesian bootstrap}
\begin{itemize}
  \item \textbf{Идея:} вместо ресэмплинга с возвращением (что эквивалентно случайным \emph{целочисленным} весам на каждой точке) генерируются непрерывные случайные веса $(w_1,\dots,w_n) \sim \mathrm{Dirichlet}(1,\dots,1)$ и используются напрямую, например $\bar X^* = \sum_i w_i X_i$.
  \item \textbf{Когда использовать:} нужна байесовская по духу оценка неопределённости без явного задания prior/likelihood; результат более гладкий, чем у классического бутстрапа (нет повторяющихся/задублированных точек).
\end{itemize}

\subsection*{Smoothed bootstrap}
\begin{itemize}
  \item \textbf{Идея:} после ресэмплинга точки $X_i^*$ добавляется небольшой шум (сглаживание ядром поверх $\widehat F_n$) вместо возврата точного исходного значения.
  \item \textbf{Когда использовать:} убирает ``комковатый''/дискретный вид сырого эмпирического распределения; полезно, когда статистика чувствительна к совпадениям значений или дискретности.
\end{itemize}

\section{Варианты бутстрапа для сложных случаев}

\subsection*{$m$-out-of-$n$ bootstrap / subsampling}
\begin{itemize}
  \item \textbf{Идея:} ресэмплируется только $m \ll n$ точек (с возвращением или без), а не $n$, результат затем масштабируется обратно.
  \item \textbf{Когда использовать:} обычный бутстрап \emph{несостоятелен} --- например, статистики на границе пространства параметров, или экстремальные статистики вроде минимума/максимума выборки (см.\ \S\ref{sec:fails}: обычный бутстрап концентрирует слишком много массы в наблюдаемом экстремуме; меньший размер ресэмпла снижает шанс каждый раз заново выбирать именно его).
  \item \textbf{На что смотреть:} нужно выбрать $m$; асимптотика требует $m\to\infty$, $m/n \to 0$.
\end{itemize}

\section{Когда бутстрап не работает}
\label{sec:fails}

Классический пример: $X_1,\dots,X_n \sim \mathrm{Uniform}[0,1]$, $M_n = \min\{X_i\}$. Истинное распределение: $nM_n \xrightarrow{d} \mathrm{Exp}(1)$ (гладкое). Бутстрап-распределение $n(M_n^*-M_n)$: с вероятностью $1-e^{-1}\approx 63\%$ бутстрап-ресэмпл ни разу заново не выбирает истинный минимум, так что $M_n^*=M_n$ точно --- огромный кусок вероятности сваливается в одну точку вместо того, чтобы гладко размазаться. \emph{(См.\ \texttt{figures/bootstrap\_failure\_plot.pdf} vs.\ \texttt{bootstrap\_success\_plot.pdf} для наглядного сравнения.)}

\textbf{Общий вывод:} бутстрап работает плохо всякий раз, когда статистика определяется небольшим числом экстремальных/граничных наблюдений (минимум, максимум, хвостовые квантили), а не гладким усреднением по всей выборке. Решение: $m$-out-of-$n$ bootstrap, либо отдельно построенная асимптотическая теория вместо ресэмплинга.

\section{Шпаргалка по применимости (подбор параметров пулов Curve)}

\begin{center}
\begin{tabular}{@{}p{3.4cm}p{5.6cm}p{5.4cm}@{}}
\toprule
\textbf{Метод} & \textbf{Где полезен} & \textbf{На что смотреть} \\
\midrule
Непараметрический (case) & CI на реализованную волатильность, impermanent loss, slippage --- без предположений о распределении & нужны i.i.d.; не работает для хвостовых/экстремальных статистик \\
Параметрический & симуляция поведения пула при заданном процессе цены до того, как оно реально произошло & мусор на входе $\Rightarrow$ мусор на выходе, если модель неверна \\
Block / stationary & ончейн-временные ряды (цены, объёмы, LP-потоки) --- зависимые, не i.i.d.\ & выбор длины блока \\
Wild & регрессия fees/slippage на состояние пула с гетероскедастичным шумом & только для остатков регрессии \\
$m$-out-of-$n$ & CI для хвостового риска / худшей просадки, риска ликвидации (экстремальные статистики) & выбор $m$ \\
\bottomrule
\end{tabular}
\end{center}

\end{document}
